\section{Limitações do Estudo}\label{sec:limitacoes}

Toda análise cienciométrica oferece um mapeamento retrospectivo valioso, mas delimitado por fronteiras metodológicas que demandam interpretação contextualizada. A primeira restrição relevante refere-se ao viés de atualidade (\emph{recency bias}): artigos inovadores publicados entre 2024 e 2026 --- abordando síntese por inteligência artificial, manufatura aditiva de ligas leves ou heterostruturas nanorrefletoras --- requerem uma janela de maturação temporal para ampla leitura e citação. Consequentemente, métricas quantitativas de impacto acumulado tendem naturalmente a privilegiar estudos pioneiros mais antigos em detrimento da fronteira mais recente.

Além disso, a relevância bibliográfica em periódicos não equivale diretamente à prontidão de voo em ambientes reais. O sucesso conceitual e a tenacidade teórica obtidos em corpos de prova de bancada laboratorial não eliminam a suscetibilidade a falhas prematuras sob os choques térmicos severos, gradientes de pressão e erosão de plasma característicos de túneis de vento hipersônicos ou tochas de plasma aerotérmicas.

Por fim, o recorte na base \emph{Web of Science Core Collection} prioriza publicações revisadas por pares em língua inglesa. Essa escolha preserva a comparabilidade internacional, mas exclui literatura cinzenta estratégica, relatórios confidenciais de defesa e agências espaciais, patentes e pesquisas publicadas exclusivamente em chinês por institutos asiáticos de vanguarda. Essa barreira linguística e de confidencialidade soma-se à heterogeneidade taxonômica persistente na literatura, em que diferentes comunidades científicas empregam termos múltiplos (HEAs, RHEAs, CCAs, MPEAs) para conceitos metalúrgicos estreitamente correlacionados.
